%%
%% Chapter 2: Main Body Example (正文示例)
%%

\chapter{非常规文本类型格式模板}

本章主要展示毕业论文中各类非常规文本类型的格式要求，包括表格、插图、公式和引用的规范格式。

\section{表格规范}

表格必须与论文叙述有直接联系，不得出现与论文叙述脱节的表格。表格中的内容不得与正文矛盾。每个表格都应有自己的标题和序号。标题应写在表格上方正中，不加标点，序号写在标题左方，表格须按章进行编序。表序必须连续，不得跳缺。

表中数据应正确无误，书写清楚。数字空缺的格内加"—"字线（占2个数字宽度）。表内文字和数字上、下或左、右相同时，不允许用"″"、"同上"之类的写法，可采用通栏处理方式。

\begin{table}[htbp]
  \centering
  \caption{合金钢的化学成分与力学性能}
  \label{tab:alloy}
  \begin{tabular}{cccccc}
    \toprule
    材料名称 & C & Mn & Cr & 其他 & 抗拉强度 $\sigma_b$（N/mm$^2$） \\
    \midrule
    合金钢A & 0.15 & 1.20 & 0.80 & — & 520 \\
    合金钢B & 0.22 & 1.50 & 1.10 & Ni:0.30 & 650 \\
    合金钢C & 0.18 & 1.00 & 0.60 & Mo:0.15 & 480 \\
    \bottomrule
  \end{tabular}
\end{table}

\section{插图规范}

插图应与文字内容相符。所有制图应符合国家标准和专业标准。对无规定符号的图形应采用该行业的常用画法。

每幅插图应有标题和序号，须按章进行编序，图序必须连续，不重复，不跳缺。插图标题写在图的下方正中。

\begin{figure}[htbp]
  \centering
  \fbox{\parbox{0.6\textwidth}{\centering\vspace{3cm}请在此处插入图片\\[3cm]}}
  \caption{示例插图标题}
  \label{fig:example}
\end{figure}

如图\ref{fig:example}所示，插图应与正文内容紧密配合。插图的坐标比例、线条粗细、字体大小等应符合出版要求。

\section{公式规范}

公式应另起一行写在中间。公式较长时最好在等号处转行。公式须按章进行编序，公式序号写在公式右侧齐右端线。

\begin{equation}
  E = mc^2
  \label{eq:einstein}
\end{equation}

如式\eqref{eq:einstein}所示，这是爱因斯坦质能方程。公式中的符号应在公式下方进行解释说明。

\begin{equation}
  \frac{\partial u}{\partial t} = \alpha \left( \frac{\partial^2 u}{\partial x^2} + \frac{\partial^2 u}{\partial y^2} \right)
  \label{eq:heat}
\end{equation}

式\eqref{eq:heat}为二维热传导方程，其中 $u$ 表示温度，$t$ 表示时间，$\alpha$ 为热扩散系数。

\section{引用规范}

引用参考文献时采用顺序编码制，即按照文章正文部分引用的先后顺序连续编码。引用时使用上标格式标注在引文右上角\upcite{ref-example-1}。

多个参考文献同时引用时，按编号从小到大排列\upcite{ref-example-2}。多次引用同一参考文献时，使用同一编号。
